% !TEX TS-program = pdflatex
% !TEX encoding = UTF-8 Unicode

% This file is a template using the "beamer" package to create slides for a talk or presentation
% - Giving a talk on some subject.
% - The talk is between 15min and 45min long.
% - Style is ornate.

% MODIFIED by Jonathan Kew, 2008-07-06
% The header comments and encoding in this file were modified for inclusion with TeXworks.
% The content is otherwise unchanged from the original distributed with the beamer package.

\documentclass{beamer}
\setbeamercovered{invisible}



% Copyright 2004 by Till Tantau <tantau@users.sourceforge.net>.
%
% In principle, this file can be redistributed and/or modified under
% the terms of the GNU Public License, version 2.
%
% However, this file is supposed to be a template to be modified
% for your own needs. For this reason, if you use this file as a
% template and not specifically distribute it as part of a another
% package/program, I grant the extra permission to freely copy and
% modify this file as you see fit and even to delete this copyright
% notice. 


\mode<presentation>
{
  \usetheme{Warsaw}
  % or ...

  % or whatever (possibly just delete it)
}


\usepackage[english]{babel}
% or whatever

\usepackage[utf8]{inputenc}
% or whatever

\usepackage{times}
\usepackage[T1]{fontenc}
% Or whatever. Note that the encoding and the font should match. If T1
% does not look nice, try deleting the line with the fontenc.

%%% MATH RELATED

\input{macros.tex}

\title{MATH 350 Advanced Calculus} 
\subtitle
{} % (optional)

\author[W.R. Casper] % (optional, use only with lots of authors)
{W.R. Casper}
% - Use the \inst{?} command only if the authors have different
%   affiliation.

\institute[California State University Fullerton] % (optional, but mostly needed)
{
  Department of Mathematics\\
  California State University Fullerton}
% - Use the \inst command only if there are several affiliations.
% - Keep it simple, no one is interested in your street address.

\subject{Talks}
% This is only inserted into the PDF information catalog. Can be left
% out. 



% If you have a file called "university-logo-filename.xxx", where xxx
% is a graphic format that can be processed by latex or pdflatex,
% resp., then you can add a logo as follows:

% \pgfdeclareimage[height=0.5cm]{university-logo}{university-logo-filename}
% \logo{\pgfuseimage{university-logo}}



% Delete this, if you do not want the table of contents to pop up at
% the beginning of each subsection:
\AtBeginSubsection[]
{
  \begin{frame}<beamer>{Outline}
    \tableofcontents[currentsection,currentsubsection]
  \end{frame}
}


% If you wish to uncover everything in a step-wise fashion, uncomment
% the following command: 

%\beamerdefaultoverlayspecification{<+->}


\begin{document}

\begin{frame}
  \titlepage
\end{frame}

\begin{frame}{Outline}
  \tableofcontents
  % You might wish to add the option [pausesections]
\end{frame}

% Since this a solution template for a generic talk, very little can
% be said about how it should be structured. However, the talk length
% of between 15min and 45min and the theme suggest that you stick to
% the following rules:  

% - Exactly two or three sections (other than the summary).
% - At *most* three subsections per section.
% - Talk about 30s to 2min per frame. So there should be between about
%   15 and 30 frames, all told.

\section{Real Analysis Lecture 2}


\subsection{Types of real numbers}

\begin{frame}{The problem with axioms}
Problems when axiomatizing:
\begin{itemize}
\pause
\item intuitive axioms
\pause
\item enough axioms
\pause
\item no redundant axioms
\pause
\item no inconsisent axioms
\end{itemize}
\end{frame}

\begin{frame}{We need to redetermine basics}
For us, we've lost some basic important things...
\begin{itemize}
\item We have $0$ and $1$ {\color{red}by axioms}
\pause
\item what is $2$?
\pause
$$\textbf{Definition:}\quad 2 = 1+1$$
\pause
\item what is $3$?
\pause
$$\textbf{Definition:}\quad 3 = 2+1$$
\pause
\item are $1$, $2$, $3$ all different?
\end{itemize}
\end{frame}

\begin{frame}{Challenge!}
\begin{prob}
Prove that $1$, $2$, and $3$ are all different.
\end{prob}
\pause
\begin{soln}
\pause
Since $0 < 1$, A7 says $0+1 < 1+1$.\\
\pause
Therefore $1<2$.\\
\pause
Then again by A7, $1+1 < 2+1$.\\
\pause
Therefore $2<3$.\\
\pause
Also transitivity (A9) says $1 < 2$ and $2 < 3$ implies $1 < 3$.\\
\pause
The trichotomy then implies $1\neq 2$ and $2\neq 3$ and $1\neq 3$.
\end{soln}
\end{frame}

\begin{frame}{Natural numbers}
\begin{quest}
What are the natural numbers $1,2,3,\dots$ based on the Axioms?
\end{quest}
\begin{defn}
\begin{itemize}
\pause
\item a set $\mathbb N$
\pause
\item $1\in \mathbb N$
\pause
\item $2\in \mathbb N$
\pause
\item $3\in \mathbb N$
\pause
\item general property:
$$x\in \mathbb N\quad\Longrightarrow\quad x+1\in \mathbb N.$$
\end{itemize}
\end{defn}
\end{frame}

\begin{frame}{Inductive sets}
\begin{defn}
A subset $S$ of $\mathbb{R}$ is called a \vocab{inductive set} if
\begin{itemize}
\pause
\item $1\in S$
\pause
\item if $x\in S$ then $x+1\in S$ 
\end{itemize}
\end{defn}
\pause
\begin{quest}
Is there more than one inductive set?
\end{quest}
\pause
$$\mathbb{R},\quad\mathbb{Q},\quad(0,\infty),\quad\mathbb{Z},\quad\mathbb{N},\dots$$
\end{frame}

\begin{frame}{Natural numbers}
\begin{itemize}
\pause
\item the set of \textbf{natural numbers} (or \textbf{positive integers}) is
$$\mathbb{N} = \{x: \text{$x$ belongs to every inductive set}\}$$
\pause
\item integers are
$$\mathbb{Z} = \{x: \text{$x$ is zero or $\pm x$ is a positive integer}\}$$
\pause
\item rationals are
$$\mathbb{Q} = \left\lbrace\frac{a}{b}: a,b\in\mathbb{Z},\ b\neq 0\right\rbrace.$$
\end{itemize}
\end{frame}

\begin{frame}{Induction}
\textbf{Principle of Induction:}
$$\text{If $S$ is an inductive set, then $\mathbb{N}\subseteq S$}$$
\pause
\begin{quest}
Waaaaait...How is this this induction?!
\end{quest}
\pause
\textbf{Example:} Let's prove $n(n+1)$ is divisible by $2$ for all positive integers $n$, using the Principle of Induction.\\
\pause
\textbf{Proof:}\\
\pause
Let
$$S = \{n\in\mathbb{N}: 2 | n(n+1)\}$$.
\pause
We see that $1\in S$ because $2$ divides $1(1+1)$.\\
\pause
Now suppose $x\in S$.  (This is our usual inductive assumption).\\
\end{frame}
\begin{frame}{Induction}
\textbf{Example:} Let's prove $n(n+1)$ is divisible by $2$ for all positive integers $n$, using the Principle of Induction.\\
\textbf{Proof continued:}\\
\pause
Then $2$ divides $x(x+1)$, so $x(x+1) = 2k$ for some integer $k$.\\
\pause
This means $(x+1)(x+2) = x(x+1) + 2(x+1) = 2(k+x+1)$.\\
\pause
Thus $2$ divides $(x+1)(x+2)$, showing that $x+1\in S$.\\
\pause
Since $x+1$ was arbitrary, this shows that $S$ is an inductive set.\\
\pause
By the Principle of Induction, this means $\mathbb{N}\subseteq S$.\\
\pause
In other words $n\in S$ for every positive integer $n$.\\
\pause
Hence $2$ divides $n(n+1)$ for every positive integers $n$.
\end{frame}

\begin{frame}{Other numbers}
\begin{quest}
What about other kinds of numbers, like irrational algebraic numbers or transcendentals?  Why do they exist {\color{red}according to the axioms}?
\end{quest}
To answer this, we need to understand the tenth axiom of the real numbers:
$$\textbf{the completeness axiom}$$
\end{frame}


\subsection{The tenth axiom}

\begin{frame}{Real numbers}
The \textbf{real numbers} are the unique field $\mathbb{R}$ satisfying Axiom 1-Axiom 9, plus an extra axiom called the completeness axiom.
$$\mathbb{R} = \text{the one and only complete, ordered field}$$
A complete ordered field $F$ satisfies
\begin{enumerate}[\text{A}1]
\setcounter{enumi}{9}
\pause
\item \textbf{completeness axiom}: every nonempty set of numbers $S\subseteq F$ which is bounded above has a supremum $\sup(S)$
\end{enumerate}
\begin{itemize}
\pause
\item Understanding what this means is challenging, but will be fundamental to the entire course!
\pause
\item Intuitively, it represents the fact that the real line has no holes or gaps.
\end{itemize}
\end{frame}
\begin{frame}{Upper bounds}
An \textbf{upper bound} for a set $S\subseteq\mathbb{R}$ is a number $b$ such that
$$x \leq b\quad\text{for all}\ x\in S.$$
\pause
In this case, we say $S$ is \textbf{bounded above} by $b$.\\
\pause
If $b\in S$ also, then $b$ is called a \textbf{maximal element} of $S$
\pause
\textbf{Examples:}
\begin{itemize}
\pause
\item $34$ is an upper bound of $[-1,5]$, but not a maximal element
\pause
\item $5$ is a maximal element of $[-1,5]$
\pause
\item $3$ is an upper bound of $[0,3)$, but not a maximal element
\pause
\item $\mathbb{N}$ has no upper bound
\pause
\item $[3,7)$ has an upper bound but no maximal element
\end{itemize}
\end{frame}


\begin{frame}{Challenge!}
\begin{prob}
Show that if $S$ has a maximal element, then it is unique.
\end{prob}
\pause
\textbf{Hint: use the definition!}
\pause
\begin{soln}
Suppose that $b_1$ and $b_2$ are both maximal elements of $S$.\\
\pause
Then $x\leq b_1$ and $x\leq b_2$ for all $x\in S$.\\
\pause
Moreover, $b_1\in S$ and $b_2\in S$.\\
\pause
Since $b_1\in S$ and $b_2$ is an upper bound of $S$, $b_1\leq b_2$.\\
\pause
Likewise, since $b_2\in S$ and $b_1$ is an upper bound of $S$, $b_2\leq b_1$.\\
\pause
By the trichotomy, we find $b_1=b_2$.
\end{soln}
\pause
Now we can say \emph{the} maximum, $\max(S)$
\end{frame}

\begin{frame}{Challenge}
\begin{prob}
Prove that if $A$ is a nonempty set of integers that is bounded above, then $A$ has a maximum.
\end{prob}
\end{frame}

\begin{frame}{Challenge}
\begin{soln}
The set $A$ is bounded above, so it has a supremum $b=\sup(A)$. \\
\pause
Since $b-1$ is not an upper bound, so there exists $a\in A$ with $b-1 < a$. \\
\pause
If $a'\in A$, then $a'$ is an integer and therefore $a' = a+k$ for $k\in \mathbb{Z}$. \\
\pause
Since $b$ is an upper bound, $b\geq a+k>b-1+k$, making $0 > k-1$. \\
\pause
Therefore $k\leq 0$ and $a'\leq a$. \\
\pause
It follows that $a$ is an upper bound of $A$, and since $a\in A$ it is a maximum.
\end{soln}
\end{frame}


\begin{frame}{Supremum}
A \textbf{supremum} of a set $S$ of real numbers is a real number $b\in\mathbb{R}$ such that
\begin{itemize}
\item $b$ is an upper bound of $S$
\item if $b'<b$, then $b'$ is not an upper bound of $S$
\end{itemize}
\pause
In other words
$$\text{a supremum is a \bf{least upper bound}}$$
\end{frame}

\begin{frame}{Challenge!}
\begin{prob}
Show that if $S$ has a supremum, then it is unique.
\end{prob}
\pause
\textbf{Hint: use the definition!}
\pause
\begin{soln}
Suppose that $b_1$ and $b_2$ are both suprema of $S$.\\
\pause
Then $b_1$ and $b_2$ are both upper bounds of $S$.\\
\pause
Since $b_1$ is a least upper bound, $b_1\leq b_2$.\\
\pause
Since $b_2$ is also a least upper bound, $b_2\leq b_1$.\\
\pause
By the trichotomy, we find $b_1=b_2$.
\end{soln}
\pause
Now we can say \emph{the} supremum, $\sup(S)$
\end{frame}

\begin{frame}{Challenge!}
\begin{prob}
Show that if $S$ has a maximum element.  Then $S$ has a supremum and $\max(S)=\sup(S)$
\end{prob}
\pause
\textbf{Hint: use the definition!}
\pause
\begin{soln}
Let $b=\max(S)$.\\
\pause
Then $b$ is an upper bound of $S$ and $b\in S$.\\
\pause
If $b'$ is another upper bound of $S$, then by definition $b\leq b'$.\\
\pause
Thus $b$ is the least upper bound of $S$.
\end{soln}
\end{frame}

\begin{frame}{Completeness Axiom}
Now we have the machinery in place to state the Completeness Axiom.
\pause
\begin{enumerate}[\text{A}1]
\setcounter{enumi}{9}
\pause
\item \textbf{completeness axiom}: if $S$ is any nonempty subset of real numbers which is bounded above, then it has a supremum $\sup(S)$
\end{enumerate}
\pause
\textbf{Examples:}
\begin{itemize}
\pause
\item $3$ is the supremum of $(0,3)$
\pause
\item $2$ is the supremum of
$$\left\lbrace
\frac{2n}{n+1}: n\in\mathbb{N}
\right\rbrace$$
\pause
\item $1$ is the supremum of
$$\left\lbrace
\sin(n): n\in\mathbb{N}
\right\rbrace$$
\end{itemize}
\end{frame}

\begin{frame}{Decimals}
The completeness axiom allows us to define \textbf{decimals}.
\pause
Let $0\leq n$ be an integer and $d_1,d_2,d_3,\dots$ be a sequence of integers with $0\leq d_k\leq 9$ for all $k\in\mathbb{N}$.  
\pause
\begin{defn}
The \textbf{$k$'th decimal apprixomation} of $n.d_0d_1d_2d_3\dots$ is
$$s_k = n+\frac{d_1}{10} + \frac{d_2}{10^2} + \frac{d_3}{10^3} + \dots + \frac{d_k}{10^k}.$$
\end{defn}
\pause
\begin{defn}
The \textbf{decimal} $n.d_0d_1d_2d_3\dots$ is the real number defined as the supremum of decimal approximations
$$n.d_0d_1d_2d_3\dots = \sup\{ s_k: k\in\mathbb N\dots\}.$$
\end{defn}
\end{frame}

\begin{frame}{Examples of decimals}
\begin{itemize}
\item Liouville's constant
$$0.1100010000000000000000010\dots$$
\pause
\item $\pi$
$$3.1415926535\dots$$
\pause
\item $1$
$$0.99999999999\dots$$
\pause
\item $1$
$$1.00000000000\dots$$
\pause
\item later we will show \emph{any} positive number has at least one decimal expansion
\end{itemize}
\end{frame}

\begin{frame}{Lower bounds and infima}
A \textbf{lower bound} for a set $S\subseteq\mathbb{R}$ is a number $b$ such that
$$b \leq x\quad\text{for all}\ x\in S.$$
\pause
In this case, we say $S$ is \textbf{bounded below} by $b$.\\
\pause
If $b\in S$ also, then $b$ is called a \textbf{minimal element} of $S$.
\pause
An \textbf{infimum} of a set $S$ of real numbers is a real number $b\in\mathbb{R}$ such that
\begin{itemize}
\item $b$ is a lower bound of $S$
\item if $b<b'$, then $b'$ is not a lower bound of $S$
\end{itemize}
\pause
In other words
$$\text{an infimum is a \bf{greatest lower bound}}$$
\end{frame}

\begin{frame}{Properties of Suprema}
The first property of suprema is that they must be arbitrarily close to elements of the set.
\begin{thm}[Approximation Property]
Let $S\subseteq\mathbb{R}$ be nonempty and bounded above, and let $b = \sup(S)$.
Then for all $\epsilon > 0$, there exists $a\in S$ with
$$b-\epsilon < a \leq b.$$
\end{thm}
\end{frame}
\begin{frame}{Properties of Suprema}
\begin{proof}
Since $b-\epsilon < b$, the definition of a supremum implies $b-\epsilon$ cannot be an upper bound.\\
\pause
Therefore thre must exist $a\in S$ with $a > b-\epsilon$.
\end{proof}
\end{frame}

\begin{frame}{Properties of Suprema}
Also, suprema play nicely with addition.
\begin{thm}[Additive Property]
Let $A,B\subseteq\mathbb{R}$ be nonempty, bounded above sets.
$$C = \{x+y: x\in A,\ y\in B\}.$$
Then $C$ is bounded above and
$$\sup(C) = \sup(A) + \sup(B).$$
\end{thm}
\end{frame}
\begin{frame}{Properties of Suprema}
\begin{proof}
Let $a = \sup(A)$, $b = \sup(B)$, and $c=\sup(C)$. \\
\pause
We will prove $c \leq a+b$ and then $a+b\leq c$. \\
\pause
First, note $a$ is an upper bound of $A$, so $x\leq a$ for all $x\in A$. \\
\pause
Also $b$ is an upper bound of $b$, so $y\leq b$ for all $y\in B$. \\
\pause
If $z\in C$, then $z=x+y$ for some $x\in A$ and $y\in B$, and therefore $z = x+y \leq a+b$. \\
\pause
Therefore $a+b$ is an upper bound of $C$. \\
\pause
This means $c\leq a+b$. \\
\pause
Next, note for all $x\in A$ and $y\leq B$ that $x\leq c-y$. \\
\pause
Therefore $c-y$  is an upper bound of $A$ and $a\leq c-y$. \\
\pause
It follows that $y\leq c-a$ for all $y\in B$. \\
\pause
Thus $c-a$ is an upper bound of $B$, and it follows $b\leq c-a$. \\
\pause
Therefore $a+b\leq c$.
\end{proof}
\end{frame}


\begin{frame}{Challenge}
\begin{prob}
Use the completeness axiom to prove that the $\mathbb N$ is not bounded above.
(Apostol Theorem 1.17)
\end{prob}
\pause
\textbf{Hint:} assume it is and consider $\sup(\mathbb N)-1$
\pause
\begin{soln}
Assume it is.\\\pause
The completeness axiom implies that $b=\sup(\mathbb N)$ exists\\\pause
The number $b$ is an upper bound and if $b' < b$, then $b'$ is not.\\\pause
Then $b-1 < b$ by Axiom 7, so $b-1$ cannot be an upper bound.\\\pause
This means there exists $n\in\mathbb N$ with $b-1 < n$.\\\pause
It follows from Axiom 7 that $b < n+1$.\\\pause
However, $n+1\in\mathbb{Z}$, so this contradicts $b$ being an upper bound.
\end{soln}
\end{frame}

\begin{frame}{Archimedian property}
\begin{thm}[Apostol Theorem 1.18]
For every $x\in\mathbb{R}$, there exists $n\in\mathbb N$ with $n>x$.
\end{thm}
\pause
\begin{proof}
If not, then $x$ is an upper bound of $\mathbb N$.
\end{proof}
\pause
\begin{thm}[Archimedian Property of Reals]
For every $x,y\in\mathbb{R}$ with $x>0$, there exists $n\in\mathbb N$ with $y<nx$.
\end{thm}
\pause
\begin{proof}
Replace $x$ with $y/x$ in the previous theorem.
\end{proof}
\end{frame}

\end{document}
